% !TeX root = complex-manifolds-and-analytic-varieties.tex

\subsection{Complex manifolds}

\begin{definition}\label{def:complex-manifold-by-altas}
	A \emph{complex manifold} of complex dimension \(n\) is a topological space \(M\) such that
	\begin{enumerate}
		\item \(M\) is Hausdorff and second countable;
		\item \(M\) is locally homeomorphic to \(\bbC^n\), i.e., for every point \(p\in M\), there exists an open neighborhood \(U\) of \(p\) and a homeomorphism \(\varphi: U\to V\subset \bbC^n\), where \(V\) is an open subset of \(\bbC^n\),
		      The pair \((U,\varphi)\) is called a \emph{chart};
		\item if \((U,\varphi)\) and \((U',\varphi')\) are two charts with \(U\cap U'\neq \emptyset\), then the transition map
		      \[
			      \varphi'\circ \varphi^{-1}: \varphi(U\cap U')\to \varphi'(U\cap U')
		      \]
		      is holomorphic.
	\end{enumerate}
	The collection of all charts \(\{(U_\alpha,\varphi_\alpha)\}\) that cover \(M\) is called an \emph{atlas}.
	If the atlas is maximal, it is called a \emph{complex structure} on \(M\).
\end{definition}

Another way to define complex manifolds is to use the language of ringed spaces.

\begin{definition}\label{def:complex-manifold-as-ringed-space}
	A \emph{complex manifold} of complex dimension \(n\) is a locally ringed space \((M,\calO_M)\) such that
	\begin{enumerate}
		\item \(M\) is Hausdorff and second countable;
		\item for every point \(p\in M\), there exists an open neighborhood \(U\) of \(p\) such that \((U,\calO_M|_U)\) is isomorphic to \((B,\calO_B)\), where \(B\) is the unit open ball in \(\bbC^n\) and \(\calO_B\) is the sheaf of holomorphic functions on \(B\).
	\end{enumerate}
\end{definition}

\begin{definition}\label{def:holomorphic_map_between_complex_manifolds}
	Let \(M\) and \(N\) be two complex manifolds.
	A continuous map \(f: M\to N\) is called \emph{holomorphic} if for every point \(p\in M\), there exist charts \((U,\varphi)\) of \(M\) around \(p\) and \((V,\psi)\) of \(N\) around \(f(p)\) with \(U \subset f^{-1}(V)\) such that
	\[
		\psi\circ f\circ \varphi^{-1}: \varphi(U)\to \psi(V)
	\]
	is holomorphic.
\end{definition}

\begin{definition}\label{def:submanifold_of_complex_manifold}
	Let \(M\) be a complex manifold of complex dimension \(n\).
	A subset \(S\subset M\) is called a \emph{complex submanifold} of complex dimension \(k\) if for every point \(p\in S\), there exist a chart \((U,\varphi)\) of \(M\) around \(p\) such that
	\[
		\varphi(U\cap S) = \varphi(U)\cap (\bbC^k\times \{0\}) \subset \bbC^n,
	\]
	where we identify \(\bbC^n\) with \(\bbC^k\times \bbC^{n-k}\).
	This gives a chart of \(S\) around \(p\).
	Endowed with the induced topology and the induced complex structure, \(S\) is a complex manifold of complex dimension \(k\).
\end{definition}

\begin{example}\label{eg:complex_vector_space_as_complex_manifold}
	Any complex vector space \(V\) of complex dimension \(n\) is a complex manifold of complex dimension \(n\).
\end{example}

\begin{example}\label{eg:complex_projective_space_as_complex_manifold}
	The complex projective space \(\bbC\bbP^n:= \bbC^{n+1}\setminus \{0\}/\bbC^\times\) is a complex manifold of complex dimension \(n\).
	In fact, \(\bbC\bbP^n\) can be covered by \(n+1\) charts, each of which is biholomorphic to \(\bbC^n\).
	For example, the chart \(U_0 = \{[z_0: z_1: \cdots: z_n]\in \bbC\bbP^n: z_0\neq 0\}\) is biholomorphic to \(\bbC^n\) via the map
	\[
		[z_0: z_1: \cdots: z_n] \mapsto \left(\frac{z_1}{z_0}, \frac{z_2}{z_0}, \ldots, \frac{z_n}{z_0}\right).
	\]
	The other charts are defined similarly.
\end{example}

\begin{proposition}\label{prop:complex_submanifold_by_holomorphic_implicit_function_theorem}
	Let \(M\) and \(N\) are complex manifolds of complex dimension \(n\) and \(m\) respectively, with \(n \geq m\).
	If \(f: M\to N\) is a holomorphic map such that \(p\) is a regular value of \(f\),
	i.e., the tangent map \(\upd f_x\) is surjective for every \(x\in f^{-1}(p)\),
	then \(f^{-1}(p)\) is a complex submanifold of \(M\) of complex dimension \(n - m\).
\end{proposition}
\begin{proof}
	For every point \(q\in f^{-1}(p)\), choose charts \((U,\varphi)\) of \(M\) around \(q\) and \((V,\psi)\) of \(N\) around \(p\) such that \(f(U) \subset V\).
	By changing coordinates if necessary, we may assume that \(\det(\partial f/\partial w)(q) \neq 0\),
	where we write the coordinates of \(\varphi(U)\) as \((z,w) = (z_1, \ldots, z_{n-m}, w_1, \ldots, w_m) \in \bbC^{n-m} \times \bbC^m\).
	Then by the Holomorphic Implicit Function Theorem (\cref{thm:holomorphic_implicit_function_theorem}),
	there exist open neighborhoods \(U'\) of \(q\) such that \(f^{-1}(p)\cap U'\) is biholomorphic to an open subset of \(\bbC^{n-m}\).
\end{proof}

\begin{example}\label{eg:non_singular_complex_algebraic_variety_as_complex_manifold}
	Let \(X \subset \bbC^n\) be a complex algebraic variety defined by the vanishing of polynomials \(f_1, \ldots, f_m \in \bbC[z_1, \ldots, z_n]\).
	Suppose that \(X\) is non-singular, i.e., for every point \(p\in X\), the Jacobian matrix \(\left( \partial_{z_j}f_i(p) \right)_{i,j}\) has maximal rank \(r\).
	Then \(X\) is a complex submanifold of \(\bbC^n\) of complex dimension \(n - r\).
\end{example}

\begin{example}\label{eg:hypersurface_in_CPn_as_complex_manifold}
	A \emph{hypersurface} \(H\) in \(\bbC\bbP^n\) is the zero locus of a homogeneous polynomial \(f \in \bbC[z_0, z_1, \ldots, z_n]\).
	Suppose \(0\) is a regular value of \(f: \bbC^{n+1}\setminus \{0\} \to \bbC\).
	On each chart \(U_i \cong \bbC^n\) of \(\bbC\bbP^n\), it defines a holomorphic function \(f_i: U_i \to \bbC, [z] \mapsto z = (z_1, \ldots, z_{i-1}, 1, z_{i+1}, \ldots, z_n) \mapsto f(z)\).
	The regularity condition implies that \(0\) is a regular value of each \(f_i\).
	Hence \(H \cap U_i = f_i^{-1}(0)\) is a complex submanifold of \(U_i\) of complex dimension \(n-1\) by \cref{prop:complex_submanifold_by_holomorphic_implicit_function_theorem}.
	Gluing these local pieces together, we see that \(H\) is a complex submanifold of \(\bbC\bbP^n\) of complex dimension \(n-1\).
	% Concretely, the Fermat hypersurface \(H\), which is defined by the equation \(z_0^m + z_1^m + \ldots + z_n^m = 0\) in \(\bbC\bbP^n\) for some integer \(m \geq 1\), is a complex manifold of complex dimension \(n-1\).
\end{example}

\begin{proposition}\label{prop:quotient_by_a_discrete_group_of_holomorphic_automorphisms}
	Let \(M\) be a complex manifold and let \(G\) be a discrete group acting on \(M\) by holomorphic automorphisms.
	If the action is free and properly discontinuous, then the quotient space \(M/G\) is a complex manifold and the quotient map \(\pi: M\to M/G\) is a holomorphic covering map.
\end{proposition}
\begin{proof}
	For every point \(p\in M/G\), choose a point \(q\in M\) such that \(\pi(q) = p\).
	Since the action is free and properly discontinuous (see \cref{rmk:properly_discontinuous_action}), there exists an open neighborhood \(U\) of \(q\) such that \(gU \cap U = \emptyset\) for all \(g\in G\setminus \{e\}\).
	Then \(\pi|_U: U \to \pi(U)\) is a homeomorphism.
	This gives a chart of \(M/G\) around \(p\).
	If we have two such charts \((\pi(U),\varphi)\) and \((\pi(U'),\varphi')\) of \(M/G\) whose intersection is non-empty, WLOG, assume that \(U \cap U' \neq \emptyset\).
	Then \(\pi^{-1}(\pi(U) \cap \pi(U')) = \bigsqcup_{g\in G} g(U \cap U')\).
	The transition map of \(U\) and \(U'\) gives the transition map of \(\pi(U)\) and \(\pi(U')\).
	Since the action of \(G\) is by holomorphic automorphisms, the transition maps are holomorphic.
\end{proof}
\begin{remark}\label{rmk:properly_discontinuous_action}
	Recall that an action of a group \(G\) on a topological space \(X\) is said to be \emph{properly discontinuous} if for every compact subset \(K \subset X\), the set \(\{g\in G: gK \cap K \neq \emptyset\}\) is finite.
	If \(G\) is a discrete group acting on a manifold \(M\) by diffeomorphisms,
	then the action is properly discontinuous and free if and only if for every point \(p\in M\),
	there exists an open neighborhood \(U\) of \(p\) such that \(gU \cap U = \emptyset\) for all \(g\in G\setminus \{e\}\).
\end{remark}

\begin{example}\label{eg:elliptic_curves_as_complex_manifolds_by_quotient_of_C}
	Let \(\Lambda \subset \bbC\) be a lattice, i.e., a discrete subgroup of \(\bbC\) generated by two \(\bbR\)-linearly independent complex numbers.
	Then \(\Lambda\) is isomorphic to \(\bbZ^2\) as an abstract group and acts on \(\bbC\) by translations, which are holomorphic automorphisms of \(\bbC\).
	Then the quotient space \(\bbC/\Lambda\) is a complex manifold of complex dimension \(1\) by \cref{prop:quotient_by_a_discrete_group_of_holomorphic_automorphisms}.
	Such a complex manifold is called an \emph{elliptic curve}.
	As real manifolds, it is diffeomorphic to \(S^1 \times S^1\).
\end{example}

\begin{example}\label{eg:hopf_manifolds_as_complex_manifolds_by_quotient_of_Cn}
	Fix \(\alpha \in \bbC^\times\) with \(|\alpha| \neq 1\).
	Let \(\bbZ\) act on \(\bbC^n\setminus \{0\}\) by \(k \cdot z = \alpha^k z\) for every \(k\in \bbZ\) and \(z\in \bbC^n\setminus \{0\}\).
	This action is free and properly discontinuous.
	Then the quotient space \((\bbC^n\setminus \{0\})/\bbZ\) is a complex manifold of complex dimension \(n\) by \cref{prop:quotient_by_a_discrete_group_of_holomorphic_automorphisms}.
	Such a complex manifold is called a \emph{Hopf manifold}.
\end{example}

\begin{example}\label{eg:Iwasawa_manifold_as_complex_manifold_by_quotient_of_complex_Heisenberg_group}
	Let
	\[ M = \left\{ \begin{pmatrix}
			1 & z_1 & z_3 \\
			0 & 1   & z_2 \\
			0 & 0   & 1
		\end{pmatrix} \;\middle|\; z_1, z_2, z_3 \in \bbC \right\} \]
	be the complex Heisenberg group, which is biholomorphic to \(\bbC^3\).
	Let \(\Gamma := M \cap \GL(3,\bbZ[\im])\).
	Then \(\Gamma\) is a discrete subgroup of \(M\) and acts on \(M\) by left multiplication, which are holomorphic automorphisms of \(M\).
	The action is free and properly discontinuous.
	Then the quotient space \(M/\Gamma\) is a complex manifold of complex dimension \(3\) by \cref{prop:quotient_by_a_discrete_group_of_holomorphic_automorphisms}.
	It is called the \emph{Iwasawa manifold}.
	One can replace \(\Gamma\) by other cocompact discrete subgroups of \(M\).
\end{example}


\subsection{Complex analytic space}

\begin{definition}\label{def:complex-analytic-space}
	A \emph{(complex) analytic space} is a locally ringed space $(X, \scrO_X)$ such that
	\begin{itemize}
		\item $X$ is a Hausdorff, locally compact and second countable;
		\item it admits an open covering by analytic sets as locally ringed space.
	\end{itemize}
\end{definition}

\begin{proposition}
	Let $X$ be a compact connected analytic space.
	Then $\scrO_X(X) = \bbC$.
\end{proposition}
\begin{proof}
	\Yang{}
\end{proof}



\begin{definition}
	\Yang{define the smooth $(p,q)$-form}
\end{definition}


\begin{theorem}\label{thm:partition-of-unity-on-analytic-space}
	Let $X$ be a complex analytic space.
	Then the sheaf $\scrO_X^{\sm,\bbR}$ admits the partition of unity.
	More explicitly, \Yang{}
\end{theorem}
\begin{proof}
	\Yang{}
\end{proof}





